% =============================================================================
% BRENT Pattern System v2 - Manuscrito para revista (Systems and Soft Computing)
% Clase Elsevier elsarticle. Compilar con:
%   pdflatex main
%   bibtex   main
%   pdflatex main
%   pdflatex main
% (o simplemente: latexmk -pdf main.tex)
%
% Este documento es AUTOCONTENIDO: no depende de ficheros de imagen externos.
% Las figuras se representan como marcadores (placeholders) con su pie, de modo
% que el PDF compila sin recursos adicionales; sustituya cada placeholder por el
% \includegraphics correspondiente cuando disponga de las figuras finales.
% =============================================================================
\documentclass[preprint,12pt,3p]{elsarticle}

\usepackage[utf8]{inputenc}
\usepackage[T1]{fontenc}
\usepackage{amsmath,amssymb}
\usepackage{booktabs}
\usepackage{graphicx}
\usepackage{xcolor}
\usepackage{multirow}
\usepackage{array}
\usepackage{siunitx}
\usepackage[hidelinks]{hyperref}
\usepackage{listings}

\journal{Systems and Soft Computing}

% --- Placeholder de figura autocontenido -------------------------------------
\newcommand{\figplaceholder}[2]{%
  \begin{figure}[!ht]
  \centering
  \fbox{\parbox[c][3.2cm][c]{0.9\linewidth}{\centering\itshape
  [Figura pendiente de insertar: #2]}}
  \caption{#2}
  \label{#1}
  \end{figure}}

\lstset{basicstyle=\ttfamily\footnotesize,breaklines=true,frame=single,
        columns=fullflexible,keepspaces=true}

\begin{document}

\begin{frontmatter}

\title{Redes neuronales convolucionales para la deteccion de anomalias en los
mercados de \emph{commodities}: un sistema de patrones del Brent debilmente
supervisado con validacion temporal robusta}

\author[a]{Autor Uno\corref{cor1}}
\ead{autor.uno@institucion.edu}
\author[b]{Autor Dos}
\cortext[cor1]{Autor de correspondencia.}
\address[a]{Afiliacion 1}
\address[b]{Afiliacion 2}

\begin{abstract}
Los mercados de \emph{commodities} estan expuestos a cambios abruptos de
regimen, incidencias de calidad de datos y patrones de precios visualmente
distintivos que resultan dificiles de captar con modelos puramente lineales.
Este articulo desarrolla un marco basado en redes neuronales convolucionales
(CNN) para el cribado de anomalias en \emph{commodities}, tomando al crudo Brent
como activo focal. El diseno empirico se apoya en un \emph{bundle} multiactivo
con 37 series de mercado efectivas y 180 variables de entrada derivadas de
\emph{commodities}, divisas, indices bursatiles y tipos de interes. Cada ventana
de 32 dias se transforma en una imagen de tres canales que combina un
\emph{Gramian Angular Summation Field} (GASF) del nivel del Brent, un
\emph{Gramian Angular Difference Field} (GADF) de los retornos recientes y un
mapa multivariante comprimido. Estas imagenes alimentan una red multitarea
EfficientNet-B1 que predice conjuntamente etiquetas de patrones chartistas
debilmente supervisadas y tres objetivos continuos (\texttt{future\_return},
\texttt{future\_low}, \texttt{future\_high}); una capa de cribado convierte las
salidas en banderas operativas de \emph{outliers}. Respecto a la version inicial
del sistema, esta segunda version (v2) incorpora una \textbf{metodologia de
validacion temporal robusta} siguiendo a Lopez de Prado: validacion cruzada
\emph{purgada} con \emph{embargo}, \emph{backtest walk-forward} y metricas de
robustez financiera (\emph{Probabilistic Sharpe Ratio}, \emph{Deflated Sharpe
Ratio} y \emph{Probability of Backtest Overfitting} via CSCV). El sistema se
empaqueta ademas como una \emph{compute capsule} reproducible y exporta todos los
resultados y la configuracion del experimento a un libro Excel para su validacion
por revisores. Los resultados deben interpretarse como una validacion funcional
del \emph{pipeline} y no como un techo de rendimiento, dado que la corrida
documentada emplea un esquema deliberadamente corto en CPU; la version v2 esta
preparada para entrenamiento prolongado en GPU.
\end{abstract}

\begin{keyword}
Redes neuronales convolucionales \sep deteccion de anomalias \sep mercados de
\emph{commodities} \sep crudo Brent \sep representacion visual de series
temporales \sep validacion \emph{walk-forward} \sep reproducibilidad
\end{keyword}

\end{frontmatter}

% =============================================================================
\section{Introduccion}
% =============================================================================
La calidad de los datos de mercado es una cuestion de primer orden para la
negociacion de \emph{commodities}, la medicion del riesgo de mercado y la
gobernanza de modelos. En entornos con modelos internos, los supervisores exigen
demostrar no solo una calibracion estadistica solida, sino tambien controles
robustos sobre la oportunidad, la completitud y la trazabilidad de los datos que
alimentan las metricas de riesgo. Este requisito es especialmente agudo en
\emph{commodities}, donde saltos abruptos, episodios de iliquidez y fricciones de
reporte generan patrones que se asemejan tanto a estres genuino de mercado como a
anomalias de datos.

La literatura ha mostrado que las CNN pueden reutilizarse fuera del dominio de
las imagenes naturales convirtiendo series temporales en objetos visuales
estructurados. Transformaciones como los \emph{Gramian Angular Fields} preservan
la geometria temporal y permiten a las CNN extraer rasgos locales y globales de
datos financieros \cite{wang2015imaging}. Las revisiones sobre \emph{deep
learning} en finanzas documentan un uso creciente de CNN para extraccion de
patrones y prediccion direccional \cite{sezer2020financial,hoseinzade2019cnnpred,
tsantekidis2020using,bao2025datadriven,giantsidi2025deep}.

Sin embargo, la mayoria de aplicaciones financieras se centran en la prediccion
direccional mas que en el cribado de anomalias, y la evidencia en
\emph{commodities} es mas limitada que en renta variable. Estudios recientes
sobre prediccion del crudo confirman la relevancia del \emph{deep learning}
\cite{shen2024intelligent,aldabagh2023hybrid,jin2024price,foroutan2024deep}, pero
el vinculo entre arquitecturas CNN basadas en imagenes y la deteccion operativa
de anomalias en Brent sigue insuficientemente explorado.

La contribucion de este trabajo es quintuple. (i) Presenta una representacion RGB
multiactivo y debilmente supervisada de las ventanas de Brent. (ii) Documenta una
arquitectura multitarea EfficientNet-B1 \cite{tan2019efficientnet} que clasifica
patrones chartistas y predice magnitudes del retorno futuro. (iii) Anade una capa
de cribado de \emph{outliers} sensible a patrones. (iv) \textbf{Novedad de la v2:}
incorpora una validacion temporal rigurosa (CV purgada con \emph{embargo},
\emph{walk-forward} y metricas PSR/DSR/PBO) siguiendo a Lopez de Prado
\cite{lopezdeprado2018advances,bailey2012sharpe,bailey2014deflated,
bailey2017probability}, corrigiendo la principal limitacion metodologica de la
version inicial. (v) Empaqueta el sistema como una \emph{compute capsule}
reproducible con exportacion integral de configuracion y resultados a Excel,
alineada con las expectativas de trazabilidad del BCE y de BCBS~239
\cite{ecb2025guide,bcbs239}.

% =============================================================================
\section{Literatura relacionada y posicionamiento}
% =============================================================================
Las raices metodologicas combinan tres corrientes. La primera es el
reconocimiento computacional de patrones de precio: Lo et al.
\cite{lo2000foundations} formalizaron el analisis estadistico de patrones
tecnicos. La segunda es la representacion visual de series temporales: Wang y
Oates \cite{wang2015imaging} mostraron que las transformaciones de imagen mejoran
la clasificacion e imputacion al mapear secuencias 1D en objetos 2D. La tercera
es la difusion del \emph{deep learning} en finanzas
\cite{sezer2020financial,hoseinzade2019cnnpred,tsantekidis2020using}.

En deteccion de anomalias, las revisiones subrayan que las anomalias en series
multivariantes son dificiles por la dependencia temporal, las interacciones
transversales, la no estacionariedad y la escasez de \emph{ground truth}
\cite{choi2021deep,li2023deep}. Los enfoques financieros suelen apoyarse en
\emph{autoencoders}, redes recurrentes o hibridos estadistico-neuronales
\cite{crepey2022anomaly,perez2024hybrid}. Los estudios de \emph{commodities} y
crudo confirman el valor de arquitecturas no lineales
\cite{shen2024intelligent,aldabagh2023hybrid,jin2024price,foroutan2024deep}, pero
optimizan el error de prediccion mas que las banderas operativas de anomalia. El
presente estudio ocupa un nicho diferenciado: Brent como activo focal dentro de un
sistema multiactivo de codificacion de imagenes cuyo resultado es un marco de
cribado de anomalias sensible a patrones.

Una diferencia central de la v2 respecto a la literatura CNN-finanzas es la
adopcion explicita del marco de validacion de Lopez de Prado
\cite{lopezdeprado2018advances}: la evaluacion solapada con \emph{stride}~1 se
sustituye por particiones purgadas con \emph{embargo} y por un \emph{backtest
walk-forward} evaluado con el \emph{Deflated Sharpe Ratio}
\cite{bailey2014deflated} y la \emph{Probability of Backtest Overfitting}
\cite{bailey2017probability}.

\begin{table}[!ht]
\centering
\caption{Literatura representativa que posiciona el marco propuesto.}
\label{tab:lit}
\small
\begin{tabular}{p{2.6cm}p{2.4cm}p{3.1cm}}
\toprule
\textbf{Estudio} & \textbf{Dominio} & \textbf{Relevancia} \\
\midrule
Lo et al. \cite{lo2000foundations} & Analisis tecnico & Base del etiquetado por patrones \\
Wang y Oates \cite{wang2015imaging} & ML series temporales & Series a imagenes (GAF) \\
Hoseinzade y Haratizadeh \cite{hoseinzade2019cnnpred} & Renta variable & CNN con contexto multivariable \\
Tsantekidis et al. \cite{tsantekidis2020using} & Microestructura & Aprendizaje de representaciones \\
Crepey et al. \cite{crepey2022anomaly} & Anomalias financieras & Casos de uso practicos \\
Shen et al. \cite{shen2024intelligent} & Prediccion crudo & Idoneidad de modelos no lineales \\
Lopez de Prado \cite{lopezdeprado2018advances} & Fin. ML & CV purgada, PSR/DSR/PBO \\
\textbf{Este estudio} & Cribado anomalias & Brent RGB + EfficientNet-B1 + reglas + validacion robusta \\
\bottomrule
\end{tabular}
\end{table}

% =============================================================================
\section{Datos, construccion del bundle y diseno experimental}
% =============================================================================
El material empirico procede del Brent Pattern System. La matriz \emph{wide raw}
contiene 7004 filas diarias y 181 columnas, de las cuales 180 son variables de
entrada utilizables y una es el \emph{target} de log-retorno \emph{forward} a un
dia. El \emph{bundle} cubre del 2 de enero de 2007 al 6 de marzo de 2026, pero las
ventanas multiactivo plenamente validas comienzan a finales de 2019 porque las
ultimas variables disponibles se asocian a ESTR y SOFR.

El \emph{pipeline} solicito originalmente 28 series de Yahoo Finance, 9 de FRED y
4 nodos AAA de la curva del BCE. El \emph{bundle} final efectivo contiene 37
series base y 180 variables. Los cuatro nodos AAA solicitados al BCE (EUR\_AAA\_2Y,
5Y, 10Y y 30Y) estan ausentes del \emph{bundle} entregado; esto se trata
explicitamente como una limitacion.

\begin{table}[!ht]
\centering
\caption{Inputs efectivos y configuracion de la corrida documentada.}
\label{tab:inputs}
\small
\begin{tabular}{p{4.6cm}p{3.4cm}}
\toprule
\textbf{Input / parametro} & \textbf{Valor} \\
\midrule
Periodo del \emph{bundle}        & 2007-01-02 a 2026-03-06 \\
Matriz \emph{raw}                & 7004 $\times$ 181 \\
Matriz de variables              & 7004 $\times$ 180 \\
Series base efectivas            & 37 \\
Variables de entrada             & 180 \\
Variable objetivo                & \texttt{BRENT\_fwd\_logret\_1} \\
Ventana \emph{lookback}          & 32 dias \\
Horizonte de retorno             & 1 dia \\
Horizonte soporte/resistencia    & 5 dias \\
Resolucion de imagen (v2 CPU)    & 96 $\times$ 96 \\
Resolucion de imagen (v2 GPU)    & 160 $\times$ 160 \\
\bottomrule
\end{tabular}
\end{table}

Las 180 variables se organizan en seis grupos: 37 cierres \emph{raw}, 37
log-retornos a un dia, 37 volatilidades moviles a 20 dias, 28 porcentajes de rango
intradia, 37 ratios de medias moviles (5/20) y 4 \emph{spreads}/ratios
(SPREAD\_WTI\_BRENT, RATIO\_WTI\_BRENT, SPREAD\_US10Y\_US2Y, BRENT\_EURUSD\_RATIO).

\figplaceholder{fig:composicion}{Composicion de las 180 variables de entrada
efectivas utilizadas en el experimento Brent.}

\begin{table}[!ht]
\centering
\caption{Particiones temporales del experimento Brent.}
\label{tab:splits}
\small
\begin{tabular}{p{2.9cm}rp{2.0cm}p{2.0cm}}
\toprule
\textbf{Particion} & \textbf{Ventanas} & \textbf{Inicio} & \textbf{Fin} \\
\midrule
Train (precomp.) & 986 & 2019-12-02 & 2024-11-20 \\
Test (precomp.)  & 247 & 2024-11-21 & 2026-03-04 \\
Train etiquetado & 887 & 2019-12-02 & 2024-05-21 \\
Validacion       & 99  & 2024-05-22 & 2024-11-20 \\
Test etiquetado  & 244 & 2024-11-21 & 2026-02-26 \\
\bottomrule
\end{tabular}
\end{table}

\paragraph{Fuga de informacion y motivacion de la v2.}
La version inicial no presenta \emph{look-ahead leakage} (cada ventana usa solo
informacion hasta su fecha final y los \emph{targets} empiezan tras el fin de la
ventana), pero emplea \emph{stride}~1, lo que crea solapamiento entre ventanas
consecutivas: la primera ventana de test comparte 31 de sus 32 observaciones con
la ultima de train. Esto no es fuga futura, pero implica que el test original no es
un \emph{backtest} purgado ni con \emph{embargo}. La v2 corrige exactamente esta
limitacion (Seccion~\ref{sec:validacion}).

% =============================================================================
\section{Metodologia}
% =============================================================================
\subsection{Etiquetado de patrones debilmente supervisado}
El \emph{bundle} aporta un unico \emph{target} numerico explicito (el log-retorno
\emph{forward} a un dia). Para la clasificacion, el sistema crea etiquetas debiles
a partir de la trayectoria del precio del Brent dentro de cada ventana de 32 dias.
Tras un suavizado por media movil de tres puntos, se identifican minimos y maximos
locales y se traducen en puntuaciones de patron. Las clases candidatas son
\texttt{double\_bottom}, \texttt{double\_top}, \texttt{ascending\_channel},
\texttt{descending\_channel}, \texttt{high\_tight\_flag}, \texttt{head\_shoulders},
\texttt{inverse\_head\_shoulders} y \texttt{range}. Una ventana recibe la clase de
mayor puntuacion si su confianza supera un umbral de 0.35; en caso contrario se
descarta del subconjunto supervisado. Solo seis clases sobreviven en la muestra
documentada, lo que refuerza el desbalance entre clases.

\subsection{Codificacion RGB de ventanas multiactivo}
Cada observacion es una ventana de $32 \times 180$. En lugar de alimentar el
tensor \emph{raw} a un modelo 1D, el sistema lo convierte en una imagen de tres
canales. El canal rojo es un GASF del nivel del Brent: si $x_t$ se reescala a
$\tilde{x}_t \in [-1,1]$, entonces $\phi_t = \arccos(\tilde{x}_t)$ y
\begin{equation}
\mathrm{GASF}_{ij} = \cos(\phi_i + \phi_j).
\end{equation}
El canal verde es un GADF de los retornos recientes,
$\mathrm{GADF}_{ij} = \sin(\phi_i - \phi_j)$. El canal azul es un mapa
multivariante comprimido obtenido escalando robustamente las 180 variables,
recortandolas a un rango estable y redimensionando a una imagen cuadrada. La
representacion es semanticamente rica pero no invertible: la red aprende un
\emph{embedding} latente, coherente con la literatura de representacion visual
\cite{wang2015imaging}.

\figplaceholder{fig:encoding}{Codificacion de tres canales: GASF de niveles, GADF
de retornos y mapa multiactivo de variables.}

\subsection{Arquitectura CNN y objetivo multitarea}
El modelo usa EfficientNet-B1 \cite{tan2019efficientnet} como \emph{backbone}
visual. Encima se anade un bloque denso compartido de 256 unidades, activacion
SiLU y \emph{dropout}, seguido de dos ramas: una cabeza de clasificacion con ocho
\emph{logits} y una cabeza de regresion con tres salidas (\texttt{future\_return},
\texttt{future\_low}, \texttt{future\_high}). El objetivo global es
\begin{equation}
\mathcal{L} = \mathcal{L}_{\text{cls}} + \lambda\, \mathcal{L}_{\text{reg}},
\qquad \lambda = 0.5,
\end{equation}
con $\mathcal{L}_{\text{cls}}$ entropia cruzada (opcionalmente ponderada por clase
o \emph{focal loss} \cite{lin2017focal} en la v2 para mitigar el desbalance) y
$\mathcal{L}_{\text{reg}}$ SmoothL1/Huber.

\subsection{Capa de cribado de \emph{outliers}}
Tras la inferencia, el sistema calcula una bandera de \emph{outlier} que combina
la confianza del modelo y diagnosticos por reglas. Una ventana se marca cuando se
cumple al menos una condicion: \texttt{weighted\_outlier\_score} $\geq 0.68$,
confianza dominante en la ventana movil $< 0.55$, o \emph{metric distance}
$\geq 2.5$. La longitud de la ventana movil es 10 y el peso acumulado objetivo es
0.99. La senal de anomalia es asi sensible a patrones pero no equivale a una
\emph{misclasificacion}.

% =============================================================================
\subsection{Validacion temporal robusta (aportacion de la v2)}
\label{sec:validacion}
% =============================================================================
La v2 sustituye la evaluacion solapada por el marco de Lopez de Prado
\cite{lopezdeprado2018advances}.

\paragraph{CV purgada con \emph{embargo}.} Dado el solapamiento inducido por el
\emph{stride}~1, las etiquetas de ventanas cercanas comparten informacion. Se
\emph{purgan} del conjunto de entrenamiento las ventanas cuyo intervalo de
informacion se solapa con el de validacion/test, y se aplica un \emph{embargo} de
$h$ observaciones tras cada bloque de validacion para eliminar la contaminacion
serial residual. En la configuracion documentada, \texttt{embargo}~$=5$.

\paragraph{\emph{Backtest walk-forward}.} La evaluacion adopta un esquema
\emph{walk-forward} de $K$ folds (ventana expansiva), reentrenando en cada paso y
concatenando las predicciones \emph{out-of-sample} en una unica trayectoria
\emph{multipath}. La estrategia direccional toma posicion segun el signo del
retorno predicho.

\paragraph{Metricas de robustez financiera.} Sobre la curva de resultados se
calculan el \emph{Probabilistic Sharpe Ratio} (PSR) \cite{bailey2012sharpe}, que
ajusta el Sharpe por asimetria y curtosis y por la longitud de la muestra; el
\emph{Deflated Sharpe Ratio} (DSR) \cite{bailey2014deflated}, que penaliza el
numero de configuraciones probadas (\emph{selection bias}); y la
\emph{Probability of Backtest Overfitting} (PBO) mediante
\emph{Combinatorially Symmetric Cross-Validation} (CSCV)
\cite{bailey2017probability}. Formalmente, el PSR estima
\begin{equation}
\widehat{\mathrm{PSR}}(\mathrm{SR}^{*}) =
\Phi\!\left( \frac{(\widehat{\mathrm{SR}} - \mathrm{SR}^{*})\sqrt{n-1}}
{\sqrt{1 - \hat{\gamma}_3 \widehat{\mathrm{SR}} +
\frac{\hat{\gamma}_4 - 1}{4}\widehat{\mathrm{SR}}^2}} \right),
\end{equation}
donde $\hat{\gamma}_3$ y $\hat{\gamma}_4$ son la asimetria y la curtosis de los
retornos y $\Phi$ es la normal estandar acumulada.

\figplaceholder{fig:walkforward}{Esquema \emph{walk-forward} con purga y
\emph{embargo}: ventanas de \emph{train}/\emph{valid}/\emph{test} desplazadas en el
tiempo.}

% =============================================================================
\section{Resultados empiricos para Brent}
% =============================================================================
La muestra etiquetada esta estructuralmente desbalanceada:
\texttt{ascending\_channel}, \texttt{descending\_channel} y \texttt{range} dominan
el subconjunto supervisado. Este desbalance se refleja en la confusion del
\emph{hold-out} y en la tendencia a sobrepredecir clases tipo canal.

\begin{table}[!ht]
\centering
\caption{Rendimiento agregado por particion (corrida corta documentada).}
\label{tab:agregado}
\scriptsize
\begin{tabular}{lrrrrrr}
\toprule
\textbf{Part.} & \textbf{n} & \textbf{Acc.\%} & \textbf{Top-2\%} &
\textbf{MAE ret.} & \textbf{Signo\%} & \textbf{R$^2$} \\
\midrule
Entren. & 887  & 85.0 & 95.3 & 3.15 & 52.20 & -1.27 \\
Valid.  & 99   & 53.5 & 77.8 & 3.17 & 47.47 & -4.59 \\
Test    & 244  & 53.3 & 81.6 & 2.64 & 52.46 & -2.43 \\
Total   & 1230 & 76.2 & 91.1 & 3.05 & 51.87 & -1.51 \\
\bottomrule
\end{tabular}
\end{table}

En el \emph{hold-out} temporal, la exactitud top-1 frente a etiquetas debiles
alcanza el 53.3\% y la top-2 el 81.6\%. Esta brecha sugiere que la representacion
capta un vecindario plausible de alternativas geometricas. La prediccion del
retorno sigue siendo dificil (MAE en test de 2.64 pp y R$^2$ negativo), coherente
con el esquema deliberadamente corto de entrenamiento.

\begin{table}[!ht]
\centering
\caption{Informe de clasificacion por clase en el conjunto de test.}
\label{tab:porclase}
\small
\begin{tabular}{lrrrr}
\toprule
\textbf{Clase} & \textbf{Prec.\%} & \textbf{Recall\%} & \textbf{F1\%} & \textbf{Sop.} \\
\midrule
ascending\_channel  & 45.6 & 81.6 & 58.5 & 76 \\
descending\_channel & 63.3 & 59.4 & 61.3 & 96 \\
range               & 61.1 & 15.3 & 24.4 & 72 \\
\midrule
Global              & --   & --   & 53.3 & 244 \\
\bottomrule
\end{tabular}
\end{table}

El modelo discrimina los regimenes tipo canal mejor que el comportamiento de
rango: \texttt{ascending\_channel} logra recall alto y precision menor, mientras
\texttt{range} tiene precision aceptable pero recall de solo 15.3\%.
Operativamente, el sistema reconoce mejor la morfologia direccional que la
congestion sin direccion.

\figplaceholder{fig:confusion}{Matriz de confusion normalizada por filas del
conjunto de test.}

La capa de anomalias permanece activa tras las cabezas de patrones y regresion. A
lo largo de las 2290 ventanas de inferencia, 942 se marcan como \emph{outliers}
(41.1\%). La mayor parte de la activacion proviene de baja confianza dominante mas
que del umbral de \emph{weighted score}: la incertidumbre es un motor principal de
las alertas.

% =============================================================================
\section{Discusion}
% =============================================================================
El experimento respalda tres conclusiones. Primera, las codificaciones
convolucionales de ventanas cortas preservan estructura util (top-2 $>$ 80\%).
Segunda, el contexto multiactivo aporta valor: combinar canales especificos del
Brent con un mapa amplio de variables evita inferir anomalias solo desde la
trayectoria del petroleo. Tercera, la capa de anomalias anade relevancia
operativa al convertir salidas brutas en alertas interpretables.

Desde el riesgo de mercado, las expectativas del BCE y de BCBS~239 enfatizan la
calidad de datos, la trazabilidad y la escalada oportuna \cite{ecb2025guide,
bcbs239}. Un cribado CNN no sustituye a la gobernanza, pero actua como control
analitico de segunda linea que prioriza ventanas para revision humana. La v2
refuerza este argumento con validacion purgada y metricas DSR/PBO, que acotan el
riesgo de \emph{overfitting} de \emph{backtest}.

Las limitaciones persisten: la corrida documentada usa un esquema corto en CPU, el
desbalance es severo, el \emph{bundle} carece de las series AAA del BCE y las
etiquetas debiles son un \emph{proxy} de la estructura de anomalia. La v2 mitiga la
limitacion de evaluacion solapada, y su empaquetado reproducible facilita la
verificacion independiente. La agenda futura incluye entrenamiento prolongado en
GPU, preentrenamiento sintetico, rebalanceo de clases, recuperacion de nodos del
BCE y restricciones monotonicas para \texttt{future\_low}/\texttt{future\_high}.

% =============================================================================
\section{Conclusion}
% =============================================================================
Este articulo consolida el proyecto Brent en un manuscrito coherente sobre
deteccion de anomalias basada en CNN en \emph{commodities}, y presenta la version
v2 del sistema, que anade validacion temporal robusta (CV purgada con
\emph{embargo}, \emph{walk-forward} y PSR/DSR/PBO) y reproducibilidad de grado
publicacion (\emph{compute capsule} y exportacion integral a Excel). Los
resultados del \emph{hold-out} son moderados en top-1 y solidos en top-2, lo que
sugiere que la representacion es informativa aunque la corrida corta no este
optimizada. Con entrenamiento mas profundo en GPU y validacion temporal estricta,
la arquitectura constituye una base creible para monitorizacion de grado
productivo.

% =============================================================================
\appendix
\section{Reproducibilidad: \emph{compute capsule} y salidas}
\label{ap:repro}
% =============================================================================
El sistema se distribuye como una \emph{compute capsule} de Code Ocean con la
estructura estandar (\texttt{code/}, \texttt{data/}, \texttt{environment/},
\texttt{metadata/}, \texttt{results/}, \texttt{run}). La ejecucion es
determinista: semillas fijadas (\texttt{seed}~$=42$, \texttt{PYTHONHASHSEED}$=0$) y
entorno \emph{pinned} por \texttt{Dockerfile} y \texttt{requirements.txt}.

\paragraph{Configuracion externa.} Todos los parametros del experimento residen en
ficheros de configuracion externos al codigo, en JSON o YAML indistintamente
(\texttt{config\_codeocean.json}, \texttt{config\_codeocean.yaml} y
\texttt{config\_gpu.yaml}). Se distingue el \emph{descriptor de la capsula}
(\texttt{metadata/metadata.yml}) de la \emph{configuracion del experimento}
(ficheros \texttt{config\_*}) y del \emph{registro de reproducibilidad}
(\texttt{results/run\_manifest.json}), que combina la configuracion resuelta con el
entorno de ejecucion (versiones, \emph{commit} de git, \emph{timestamp}).

\paragraph{Salida en Excel para validacion.} El \emph{pipeline} exporta un libro
Excel (\texttt{results/reports/brent\_resultados.xlsx}) con hojas dedicadas a:
resumen, \textbf{configuracion completa del experimento}, entorno, metricas de
clasificacion (globales y por clase), matriz de confusion, regresion,
\emph{backtest} (Sharpe, PSR, DSR, PBO), Sharpe por \emph{fold}, resumen de
\emph{folds} y predicciones \emph{out-of-sample}. Esta trazabilidad permite a una
revista reproducir y validar cada cifra frente a la configuracion exacta que la
produjo.

\paragraph{Ejecucion en GPU.} La configuracion \texttt{config\_gpu.yaml} activa
\emph{mixed precision} (AMP), \emph{cuDNN benchmark}, \emph{caching} de imagenes,
mayor resolucion y regimen de entrenamiento prolongado, manteniendo identica la
metodologia respecto a la corrida CPU reproducible.

% =============================================================================
\section*{Declaraciones}
\textbf{Financiacion:} [A completar por los autores].\\
\textbf{Conflictos de interes:} [A completar por los autores].\\
\textbf{Disponibilidad de datos:} los resultados se basan en artefactos del
proyecto Brent suministrados por los autores; las series brutas proceden de Yahoo
Finance, FRED y el BCE.

% =============================================================================
\bibliographystyle{elsarticle-num}
\bibliography{references}

\end{document}
